\documentclass[a4paper,11pt]{article}
\usepackage[margin=27mm]{geometry}
\usepackage{fontspec,amsmath,amssymb,amsthm,unicode-math,booktabs}
\setmainfont{STIXGeneral}\setmathfont{STIX Two Math}
\setlength{\emergencystretch}{2em}
\newcommand{\gen}[1]{\langle#1\rangle}
\newtheorem{theorem}{Theorem}
\title{Splitting a middle generator in the symmetric group}
\author{Human Compression project}\date{Private draft; independent review pending}
\begin{document}\maketitle
\begin{abstract}
We apply the common supported-transposition theorem to the complete
principal rank-augmentation statement of De Saedeleer--Leemans--Mulpas,
in its specified author edition. The original four- and five-point cases
are proved with their original hypotheses. The six-point assertion is
false as printed; we give an exact replacement and bind a complete
counterexample. We also verify all three finite examples used by the
source to motivate its additional assumptions. The introductory census
is retained as background data rather than asserted as a new enumeration.
\end{abstract}

\section{The original statement and the corrected conclusion}
The source is the nineteen-page author edition
arXiv:2212.12563v1, 23 December 2022 [DLM]. Its Theorem~4.1 starts on
physical page~5 and continues on page~6. We retain its distinction
between the natural permutation group $S_n$ and the particular chosen
marking. Write its original rank-three C-marking as $(a,q,c)$.
Choose a transposition $t$ from the disjoint cycle expression of $q$,
with neither endpoint moved by $c$, and write $q=tb$. The proposed
new marking is $(t,a,b,c)$. Put
\[
 B=\gen{a,b},\quad H=\gen{t,a,b},\quad K=\gen{a,b,c},
\]
and let $O$ be the $H$-orbit containing the endpoints of $t$.

\begin{theorem}
Assume the complete hypotheses of [DLM, Theorem~4.1]: $n\geq5$;
$(a,q,c)$ is a rank-three C-marking of the natural $S_n$; $q$ is not a
transposition; the chosen $q$-edge is not adjacent to a $c$-edge;
and every $B$-orbit has size at most three. Then the following hold.
\begin{enumerate}
\item If $|O|=4$ and $K$ acts imprimitively on each of its orbits
containing more than one $B$-triple, the new marking is C.
\item If $|O|=5$, $a,b,c$ are even, and the same imprimitivity
condition holds, the new marking is C.
\item If $|O|=6$, let $N$ be the group generated by the three matching
transpositions between its two marked $B$-triples. The new marking is C
if and only if $N\cap K=1$. It is C, in particular, when those two triples
belong to distinct $K$-orbits. Without an additional condition, it need
not be C.
\end{enumerate}
\end{theorem}
\begin{proof}
The new marking generates the natural $S_n$, since it generates the
original marks $a,q=tb,c$. The chosen transposition commutes with both
$b$ and $c$, and the old string condition gives $ac=ca$.
The local component argument in the common paper gives exactly its
four-, five- and six-point paths; in particular the four marks are distinct
nonidentity involutions. Its two-end lemma therefore applies.
The common block-and-parity theorem proves (1) and (2). Its theorem on the six-point action proves the exact criterion and sufficient condition in (3).
The frozen eighteen-point counterexample, recalled below, disproves the
unrestricted assertion. Thus the result preserves both valid source
cases and replaces its false third case explicitly.
\end{proof}

The common proof also establishes the elementary dihedral description
in source Lemma~5.2, both end C-markings in Sections~5.1--5.2, and both
conclusions of its block Lemma~5.3. In particular, no failure of the last
intersection is being inferred merely from an error in the source's
proof: the next section gives a group element in the wrong intersection.

\section{The actual six-point obstruction}
On eighteen points use
\begin{align*}
 a&=(1\,2)(5\,6)(7\,8)(11\,12)(13\,14)(17\,18),\\
 b&=(2\,3)(4\,5)(7\,9)(10\,11)(14\,15)(16\,17),\\
 c&=(3\,4)(5\,7)(6\,8)(9\,10)(11\,13)(12\,14)(15\,16),\\
 t&=(1\,18),\qquad q=tb.
\end{align*}
The complete written proof in the separately frozen correction establishes
that $(a,q,c)$ is an irreducible C-marking of the natural $S_{18}$, of
type $[6,90]$. It supplies the whole permutations, a seventeen-edge
conjugate-transposition tree proving generation, and the exact dihedral
intersection proving all original parabolic identities. Its input
hypotheses are part of this application's full dependency boundary, not
an unpriced citation to a short diagnostic.

The six $B$-orbits are consecutive triples. The $H$-orbit containing
$t$ is $\{1,2,3,16,17,18\}$, of size six. With products acting on the
right, put $w=a(bc)^2ba(cb)^8a$. Direct multiplication gives
\[
 (bc)^{45}=(6\,8),\qquad 6^w=1,\quad8^w=18,
 \qquad ((bc)^{45})^w=t.
\]
Thus $t\in K$ and $N\cap K\ne1$. The label sets $\{t\}$ and
$\{a,b,c\}$ are disjoint, so this is an actual failure of the
intersection property. All thirteen statements in the separate Lean
supplement concern these literal output permutations and generated
subgroups. They do not formalize the old C-marking, natural symmetric
group identification, or source theorem hypotheses, which retain their
complete written proofs.

\section{The three motivating examples}
Section~6 of [DLM], physical pages~17--18, supplies examples showing
why the additional assumptions in the first two cases matter. We keep
its exact point labels. The original triples $(a,q,c)$ and the split
transpositions are as follows:
\begin{align*}
n=9:\quad
 a&=(2\,3)(6\,9)(7\,8),&
 q&=(1\,2)(3\,4)(5\,6),\\
 c&=(4\,5)(6\,7)(8\,9),&t&=(1\,2);\\[2pt]
n=10:\quad
 a&=(1\,2)(3\,4)(9\,10),&
 q&=(2\,3)(4\,5)(6\,7)(8\,9),\\
 c&=(5\,6)(7\,8),&t&=(2\,3);\\[2pt]
n=13:\quad
 a&=(1\,2)(3\,4)(7\,9)(8\,10),&
 q&=(2\,3)(4\,5)(6\,7)(10\,11)(12\,13),\\
 c&=(5\,6)(7\,8)(9\,10)(11\,12),&t&=(2\,3).
\end{align*}
In each row $b=tq$. The complete required file \texttt{check\_examples.g}
constructs these actual permutation groups. It verifies that the original
group has order $n!$, checks its involution and string conditions, and
tests every one of its sixty-four parabolic pairs. It also computes the
actual subgroup intersection and the claimed permutation action. The
resulting orders are
\[
\begin{array}{c|rrrr|c}
 n&|B|&|H|&|K|&|H\cap K|&\text{signs of }(a,b,c)\\\hline
9&12&288&40320&72&(-,+,-)\\
10&12&1440&768&24&(-,-,+)\\
13&12&1440&39916800&72&(+,+,+).
\end{array}
\]
These are finite mathematical premises for these three examples only;
the common quantified proofs do not use sample computations.

For $n=9$, $K$ fixes $1$ and acts as $S_8$ on $\{2,\ldots,9\}$.
It is primitive there. The element $(2\,3\,4)$ belongs to $H\cap K$
but not to $B$. Thus the four-point case can fail without its
imprimitivity condition.

For $n=10$, the action of $K$ on $\{3,\ldots,10\}$ preserves the
nontrivial block system
\[
 \{3,10\},\quad\{4,9\},\quad\{5,8\},\quad\{6,7\}.
\]
The program verifies that this is exactly the orbit of the first displayed
pair under $K$. But $a,b$ are odd, and the pure transposition $(1\,2)$
is in $H\cap K$ and not in $B$. Thus the parity condition cannot simply
be omitted from the five-point assertion.

For $n=13$, all three marks $a,b,c$ are even, but the projection of $K$
to $\{3,\ldots,13\}$ is the primitive natural $S_{11}$ action.
The even cycle $(3\,4\,5)$ is in $H\cap K$ but not in $B$.
Here the remaining assumption that fails is imprimitivity.
The full permutation group $K$ is consistent with evenness: an odd
permutation in that eleven-point projection has the compensating
transposition on the other pair of points. These verified examples
recover the source's motivating conclusions without identifying
projection parity with parity of the whole permutation.

\section{Scope, retained inputs and verification}
The operation and valid cases are due to De Saedeleer, Leemans and
Mulpas. We use the full common paper, the complete eighteen-point
correction proof, and the entire required example program. All are
included in the proof boundary. No cost of common theory is silently
borrowed from a different topic. The standard McMullen--Schulte end
criterion is explicitly retained. The source's quoted rank-reduction
result is background, already treated in the separate corrected L037
application; it is not needed as a premise of the positive augmentation
proof. Neither is its quoted sesqui-extension result.

The closing sentence on source physical page~18 also needs a bond-index
correction. On the distinguished $t,a$-orbit, $ta$ is a three-cycle in
the four-point row and a four-cycle in the other rows; on its complement
it acts as $a$, of order at most two. Hence the first bond $|ta|$ is three
or six in the four-point row, and four in the five- and six-point rows.
The orbit bound always makes the \emph{second} bond $|ab|$ three or six,
as proved by the common dihedral calculation. It does not force the
first bond to have those two values.

The numerical Table~1 on source physical page~2 reports a motivating
computer experiment on $S_5,\ldots,S_{11}$. Those census entries are
retained background data, not regenerated or claimed as a new result
here. The broader extensions suggested in its conclusion, to larger
middle orbits and to higher ranks, remain outside this theorem.
The generic bounded diagnostic \texttt{check\_models.g} tests small
instances of the common theorem, including actual subgroup identities;
it is a validation file, not a proof of the unbounded statements.
There is no new conjecture or numerical compression claim in this draft.

\begin{thebibliography}{9}
\bibitem{DLM} J. De Saedeleer, D. Leemans and J. Mulpas,
\emph{A rank augmentation theorem for rank 3 string C-group representations
of the symmetric groups}, arXiv:2212.12563v1, 23 December 2022.
\end{thebibliography}
\end{document}
